% TOP_PROBLEM: {"id": 3074, "title": "Big Line or Big Clique in Planar Point Sets", "category_id": 6, "category": {"id": 6, "name": "geometry", "display_name": "Geometry", "description": "Euclidean and non-Euclidean geometry, geometric structures.", "slug": "geometry", "order_index": 6, "created_at": "2026-07-31T15:26:25.671Z"}, "status": "open", "problem_number": "OPG-588", "set_id": 12}
% FIRST_POSTED: 2026-10-01
% ENTRY_KIND: statement_only
% UnsolvedMath ID 3074; source catalogue code OPG-588
% !TeX program = lualatex
% Definition and sources only; this file is not an LLM solution attempt.
\documentclass[11pt,a4paper]{article}
\usepackage[margin=25mm]{geometry}
\usepackage{fontspec}
\setmainfont{lmroman10-regular.otf}[BoldFont=lmroman10-bold.otf,
 ItalicFont=lmroman10-italic.otf,BoldItalicFont=lmroman10-bolditalic.otf]
\usepackage{amsmath,amssymb,mathtools}
\usepackage{unicode-math}
\setmathfont{latinmodern-math.otf}
\AtBeginDocument{\let\setminus\smallsetminus}
\usepackage{xurl,hyperref}
\hypersetup{colorlinks=true,urlcolor=blue}
\setcounter{secnumdepth}{0}
\setlength{\parindent}{0pt}
\setlength{\parskip}{5pt}
\setlength{\emergencystretch}{3em}
\begin{document}
\section*{Big Line or Big Clique in Planar Point Sets}\label{3074}
{\small UnsolvedMath ID 3074 · OPG-588 · Geometry · OpenGarden}
\subsection{Definitions and mathematical statement}
Let $S\subset\mathbb R^2$ be a finite set of distinct points. For
$x,y\in S$ with $x\ne y$, their open segment is
$(x,y)=\{(1-t)x+ty:0<t<1\}$. Say that $x$ and $y$ are visible with
respect to $S$ if $(x,y)\cap S=\varnothing$. The visibility graph has
vertex set $S$ and joins precisely the visible pairs.

The conjecture is that for all integers $k,\ell\ge2$, there exists an
integer $N(k,\ell)$ such that every finite $S$ with $|S|\ge N(k,\ell)$
satisfies at least one of the following alternatives:
\[
 \exists\text{ line }L:\ |S\cap L|\ge\ell,
 \qquad\text{or}\qquad
 \exists T\subseteq S:\ |T|=k\text{ and every pair in }T\text{ is visible}.
\]
Thus the second alternative is a clique of order $k$ in the visibility
graph. The threshold may depend on the two requested sizes, but not on
the coordinates of $S$.

Visibility is tested against the entire original point set, including
points outside $T$. Merely choosing $k$ points with no three collinear
does not ensure visibility: other points of $S$ may obstruct their
segments. Conversely, pairwise visibility does not require the convex hull
of $T$ to contain no additional points away from those segments.
Collinear configurations are allowed, and there is no general-position
assumption. The assertion is specifically about finite sets of sufficiently
large cardinality; replacing them by an infinite set is a different statement.
\subsection{Short English statement}
Is every sufficiently large finite planar point set forced to contain many collinear points or a large set of pairwise visible points?
\subsection{Sources}
\begin{itemize}
\item[S1] ULAM.AI and the UnsolvedMath Contributors, supplied problems.json, record 3074 (OPG-588), OpenGarden collection. Catalogue identity and source transcription; CC BY 4.0.
\url{https://huggingface.co/datasets/ulamai/UnsolvedMath}.
\item[S2] Open Problem Garden, Big Line or Big Clique in Planar Point Sets. Original problem and discussion; the mathematical formulation below is expanded from the supplied source record.
\url{http://www.openproblemgarden.org/op/big_line_or_big_clique_in_planar_point_sets}.
\end{itemize}
\end{document}
